% concatenated from UnitaryPreservation.tex
\def\CTeXPreproc{Created by ctex v0.2.12, don't edit!}









\documentclass[twoside]{amsart}
\usepackage{mathrsfs}
\usepackage{amsfonts}
\usepackage{amssymb}
\usepackage{amssymb}
\usepackage{amssymb}
\usepackage{amssymb}
\usepackage{amssymb}
\usepackage{amssymb}
\usepackage{amssymb}
\usepackage{amssymb}
\usepackage{amsmath}
\usepackage{mathdots}
\usepackage{xcolor}
\usepackage[all,cmtip]{xy}
%\usepackage{diagrams}
%\usepackege[all]{xy}
%\usepackage[small,nohug,heads=vee]{diagrams}
%\diagramstyle[labelstyle=\scriptstyle]


\usepackage{lastpage}



\RequirePackage{amsmath} \RequirePackage{amssymb}
\usepackage{amscd,latexsym,amsthm,amsfonts,amssymb,amsmath,amsxtra}
\usepackage[colorlinks=true, urlcolor=blue, linkcolor=red,anchorcolor=blue,citecolor=blue]{hyper ref}

%\usepackage[colorlinks=true, urlcolor=blue,bookmarks=true,bookmarksopen=true,
%citecolor=blue,hypertex]{hyperref}
%cohomology groups
\newcommand{\Hf}{H_{f}}
\newcommand{\HfF}{H_{f,\F}}
\newcommand{\HfFa}{H_{f,\F_{\a}}}
\newcommand{\HfFda}{H_{f,\F_{\a}^{*}}}
\newcommand{\HfFd}{H_{f,\F^{*}}}
\newcommand{\Hs}{H_{s}}
\newcommand{\HsF}{H_{s,\F}}
\newcommand{\HsFd}{H_{s,\F^{*}}}


% newcommand bb
    \newcommand{\BA}{{\mathbb {A}}} \newcommand{\BB}{{\mathbb {B}}}
    \newcommand{\BC}{{\mathbb {C}}} \newcommand{\BD}{{\mathbb {D}}}
    \newcommand{\BE}{{\mathbb {E}}} \newcommand{\BF}{{\mathbb {F}}}
    \newcommand{\BG}{{\mathbb {G}}} \newcommand{\BH}{{\mathbb {H}}}
    \newcommand{\BI}{{\mathbb {I}}} \newcommand{\BJ}{{\mathbb {J}}}
    \newcommand{\BK}{{\mathbb {K}}} \newcommand{\BL}{{\mathbb {L}}}
    \newcommand{\BM}{{\mathbb {M}}} \newcommand{\BN}{{\mathbb {N}}}
    \newcommand{\BO}{{\mathbb {O}}} \newcommand{\BP}{{\mathbb {P}}}
    \newcommand{\BQ}{{\mathbb {Q}}} \newcommand{\BR}{{\mathbb {R}}}
    \newcommand{\BS}{{\mathbb {S}}} \newcommand{\BT}{{\mathbb {T}}}
    \newcommand{\BU}{{\mathbb {U}}} \newcommand{\BV}{{\mathbb {V}}}
    \newcommand{\BW}{{\mathbb {W}}} \newcommand{\BX}{{\mathbb {X}}}
    \newcommand{\BY}{{\mathbb {Y}}} \newcommand{\BZ}{{\mathbb {Z}}}
    %newcommand textbf
      \newcommand{\bA}{{\textbf {A}}}\newcommand{\bB}{{\textbf {B}}}  \newcommand{\bC}{{\textbf {C}}}\newcommand{\bD}{{\textbf {D}}}
   \newcommand{\bE}{{\textbf {E}}}\newcommand{\bF}{{\textbf {F}}}  \newcommand{\bG}{{\textbf {G}}}\newcommand{\bH}{{\textbf {H}}}
    \newcommand{\bI}{{\textbf {I}}}\newcommand{\bJ}{{\textbf {J}}}  \newcommand{\bK}{{\textbf {K}}}\newcommand{\bL}{{\textbf {L}}}
   \newcommand{\bM}{{\textbf {M}}}\newcommand{\bN}{{\textbf {N}}}  \newcommand{\bO}{{\textbf {O}}}
    \newcommand{\bP}{{\textbf {P}}}\newcommand{\bQ}{{\textbf {Q}}}  \newcommand{\bR}{{\textbf {R}}}\newcommand{\bS}{{\textbf {S}}}
   \newcommand{\bT}{{\textbf {T}}}\newcommand{\bU}{{\textbf {U}}}  \newcommand{\bV}{{\textbf {V}}}\newcommand{\bW}{{\textbf {W}}}
    \newcommand{\bX}{{\textbf {X}}}\newcommand{\bY}{{\textbf {Y}}}\newcommand{\bZ}{{\textbf {Z}}}





% newcommand  scr
    \newcommand{\sA}{{\mathscr {A}}} \newcommand{\sB}{{\mathscr {B}}}
    \newcommand{\sC}{{\mathscr {C}}} \newcommand{\sD}{{\mathscr {D}}}
    \newcommand{\sE}{{\mathscr {E}}} \newcommand{\sF}{{\mathscr {F}}}
    \newcommand{\sG}{{\mathscr {G}}} \newcommand{\sH}{{\mathscr {H}}}
    \newcommand{\sI}{{\mathscr {I}}} \newcommand{\sJ}{{\mathscr {J}}}
    \newcommand{\sK}{{\mathscr {K}}} \newcommand{\sL}{{\mathscr {L}}}
    \newcommand{\sN}{{\mathscr {N}}} \newcommand{\sM}{{\mathscr {M}}}
    \newcommand{\sO}{{\mathscr {O}}} \newcommand{\sP}{{\mathscr {P}}}
    \newcommand{\sQ}{{\mathscr {Q}}} \newcommand{\sR}{{\mathscr {R}}}
    \newcommand{\sS}{{\mathscr {S}}} \newcommand{\sT}{{\mathscr {T}}}
    \newcommand{\sU}{{\mathscr {U}}} \newcommand{\sV}{{\mathscr {V}}}
    \newcommand{\sW}{{\mathscr {W}}} \newcommand{\sX}{{\mathscr {X}}}
    \newcommand{\sY}{{\mathscr {Y}}} \newcommand{\sZ}{{\mathscr {Z}}}


% newcommand cal
    \newcommand{\CA}{{\mathcal {A}}} \newcommand{\CB}{{\mathcal {B}}}
    \newcommand{\CC}{{\mathcal {C}}} \renewcommand{\CD}{{\mathcal {D}}}
    \newcommand{\CE}{{\mathcal {E}}} \newcommand{\CF}{{\mathcal {F}}}
    \newcommand{\CG}{{\mathcal {G}}} \newcommand{\CH}{{\mathcal {H}}}
    \newcommand{\CI}{{\mathcal {I}}} \newcommand{\CJ}{{\mathcal {J}}}
    \newcommand{\CK}{{\mathcal {K}}} \newcommand{\CL}{{\mathcal {L}}}
    \newcommand{\CM}{{\mathcal {M}}} \newcommand{\CN}{{\mathcal {N}}}
    \newcommand{\CO}{{\mathcal {O}}} \newcommand{\CP}{{\mathcal {P}}}
    \newcommand{\CQ}{{\mathcal {Q}}} \newcommand{\CR}{{\mathcal {R}}}
    \newcommand{\CS}{{\mathcal {S}}} \newcommand{\CT}{{\mathcal {T}}}
    \newcommand{\CU}{{\mathcal {U}}} \newcommand{\CV}{{\mathcal {V}}}
    \newcommand{\CW}{{\mathcal {W}}} \newcommand{\CX}{{\mathcal {X}}}
    \newcommand{\CY}{{\mathcal {Y}}} \newcommand{\CZ}{{\mathcal {Z}}}

    % newcommand frak
     \newcommand{\fa}{{\mathfrak{a}}}  \newcommand{\fb}{{\mathfrak{b}}}
     \newcommand{\fc}{{\mathfrak{c}}}  \newcommand{\fd}{{\mathfrak{d}}}
     \newcommand{\fe}{{\mathfrak{e}}}  \newcommand{\ff}{{\mathfrak{f}}}
     \newcommand{\fg}{{\mathfrak{g}}}  \newcommand{\fh}{{\mathfrak{h}}}
     \newcommand{\fii}{{\mathfrak{i}}}  \newcommand{\fj}{{\mathfrak{j}}}
     \newcommand{\fk}{{\mathfrak{k}}}  \newcommand{\fl}{{\mathfrak{l}}}
     \newcommand{\fm}{{\mathfrak{m}}}  \newcommand{\fn}{{\mathfrak{n}}}
     \newcommand{\fo}{{\mathfrak{o}}}  \newcommand{\fp}{{\mathfrak{p}}}
     \newcommand{\fq}{{\mathfrak{q}}}  \newcommand{\fr}{{\mathfrak{r}}}
     \newcommand{\fs}{{\mathfrak{s}}}  \newcommand{\ft}{{\mathfrak{t}}}
     \newcommand{\fu}{{\mathfrak{u}}}  \newcommand{\fv}{{\mathfrak{v}}}
     \newcommand{\fw}{{\mathfrak{w}}}  \newcommand{\fx}{{\mathfrak{x}}}
     \newcommand{\fy}{{\mathfrak{y}}}  \newcommand{\fz}{{\mathfrak{z}}}

    \newcommand{\fA}{{\mathfrak{A}}}  \newcommand{\fB}{{\mathfrak{B}}}
     \newcommand{\fC}{{\mathfrak{C}}}  \newcommand{\fD}{{\mathfrak{D}}}
     \newcommand{\fE}{{\mathfrak{E}}}  \newcommand{\fF}{{\mathfrak{F}}}
     \newcommand{\fG}{{\mathfrak{G}}}  \newcommand{\fH}{{\mathfrak{H}}}
     \newcommand{\fI}{{\mathfrak{I}}}  \newcommand{\fJ}{{\mathfrak{J}}}
     \newcommand{\fK}{{\mathfrak{K}}}  \newcommand{\fL}{{\mathfrak{L}}}
     \newcommand{\fM}{{\mathfrak{M}}}  \newcommand{\fN}{{\mathfrak{N}}}
     \newcommand{\fO}{{\mathfrak{O}}}  \newcommand{\fP}{{\mathfrak{P}}}
     \newcommand{\fQ}{{\mathfrak{Q}}}  \newcommand{\fR}{{\mathfrak{R}}}
     \newcommand{\fS}{{\mathfrak{S}}}  \newcommand{\fT}{{\mathfrak{T}}}
     \newcommand{\fU}{{\mathfrak{U}}}  \newcommand{\fV}{{\mathfrak{V}}}
     \newcommand{\fW}{{\mathfrak{W}}}  \newcommand{\fX}{{\mathfrak{X}}}
     \newcommand{\fY}{{\mathfrak{Y}}}  \newcommand{\fZ}{{\mathfrak{Z}}}



 % newcommand :rm
     \newcommand{\RA}{{\mathrm {A}}} \newcommand{\RB}{{\mathrm {B}}}
    \newcommand{\RC}{{\mathrm {C}}} \newcommand{\RD}{{\mathrm {D}}}
    \newcommand{\RE}{{\mathrm {E}}} \newcommand{\RF}{{\mathrm {F}}}
    \newcommand{\RG}{{\mathrm {G}}} \newcommand{\RH}{{\mathrm {H}}}
    \newcommand{\RI}{{\mathrm {I}}} \newcommand{\RJ}{{\mathrm {J}}}
    \newcommand{\RK}{{\mathrm {K}}} \newcommand{\RL}{{\mathrm {L}}}
    \newcommand{\RM}{{\mathrm {M}}} \newcommand{\RN}{{\mathrm {N}}}
    \newcommand{\RO}{{\mathrm {O}}} \newcommand{\RP}{{\mathrm {P}}}
    \newcommand{\RQ}{{\mathrm {Q}}} \newcommand{\RR}{{\mathrm {R}}}
    \newcommand{\RS}{{\mathrm {S}}} \newcommand{\RT}{{\mathrm {T}}}
    \newcommand{\RU}{{\mathrm {U}}} \newcommand{\RV}{{\mathrm {V}}}
    \newcommand{\RW}{{\mathrm {W}}} \newcommand{\RX}{{\mathrm {X}}}
    \newcommand{\RY}{{\mathrm {Y}}} \newcommand{\RZ}{{\mathrm {Z}}}
    \newcommand{\diag}{{\mathrm {diag}}}
     \newcommand{\rd}{{\mathrm {d}}}

    \newcommand{\ri}{{\mathrm {i}}}\newcommand{\Sch}{{\mathbf {sch}}}
    \newcommand{\Sets}{{\mathbf {Sets}}} \newcommand{\Nm}{{\mathrm {Nm}}}
    \newcommand{\Diff}{{\mathrm {Diff}}}\newcommand{\lenth}{{\mathrm {\lenth}}}

    \newcommand{\Ad}{{\mathrm{Ad}}} \newcommand{\Aut}{{\mathrm{Aut}}}
    \newcommand{\Br}{{\mathrm{Br}}} \newcommand{\Ch}{{\mathrm{Ch}}}
    \newcommand{\cod}{{\mathrm{cod}}} \newcommand{\cont}{{\mathrm{cont}}}
    \newcommand{\cl}{{\mathrm{cl}}}   \newcommand{\Cl}{{\mathrm{Cl}}}
    \newcommand{\disc}{{\mathrm{disc}}}\newcommand{\Eis}{{\mathrm{Eis}}}
    \newcommand{\Div}{{\mathrm{Div}}} \renewcommand{\div}{{\mathrm{div}}}
    \newcommand{\End}{{\mathrm{End}}} \newcommand{\Frob}{{\mathrm{Frob}}}
    \newcommand{\Gal}{{\mathrm{Gal}}} \newcommand{\GL}{{\mathrm{GL}}}
    \newcommand{\Hom}{{\mathrm{Hom}}} \renewcommand{\Im}{{\mathrm{Im}}}
    \newcommand{\Ind}{{\mathrm{Ind}}} \newcommand{\ind}{{\mathrm{ind}}}
    \newcommand{\inv}{{\mathrm{inv}}}
    \newcommand{\Isom}{{\mathrm{Isom}}} \newcommand{\Jac}{{\mathrm{Jac}}}
    \newcommand{\ad}{{\mathrm{ad}}}  \newcommand{\Tr}{{\mathrm{Tr}}}
    \newcommand{\Ker}{{\mathrm{Ker}}} \newcommand{\Ros}{{\mathrm{Ros}}}
    \newcommand{\Lie}{{\mathrm{Lie}}} \newcommand{\Hol}{{\mathrm{Hol}}}

    \newcommand{\cyc}{{\mathrm{cyc}}}\newcommand{\id}{{\mathrm{id}}}
    \newcommand{\new}{{\mathrm{new}}} \newcommand{\NS}{{\mathrm{NS}}}
    \newcommand{\ord}{{\mathrm{ord}}} \newcommand{\rank}{{\mathrm{rank}}}
    \newcommand{\PGL}{{\mathrm{PGL}}} \newcommand{\Pic}{\mathrm{Pic}}
    \newcommand{\cond}{\mathrm{cond}} \newcommand{\Is}{{\mathrm{Is}}}
    \renewcommand{\Re}{{\mathrm{Re}}} \newcommand{\reg}{{\mathrm{reg}}}
    \newcommand{\Res}{{\mathrm{Res}}} \newcommand{\Sel}{{\mathrm{Sel}}}
    \newcommand{\RTr}{{\mathrm{Tr}}} \newcommand{\alg}{{\mathrm{alg}}}
    \newcommand{\PSL}{{\mathrm{PSL}}}


\newcommand{\coker}{{\mathrm{coker}}}
\newcommand{\val}{{\mathrm{val}}} \newcommand{\sign}{{\mathrm{sign}}}
\newcommand{\mult}{{\mathrm{mult}}} \newcommand{\Vol}{{\mathrm{Vol}}}
\newcommand{\Meas}{{\mathrm{Meas}}}\renewcommand{\mod}{\ \mathrm{mod}\ }
\newcommand{\Ann}{\mathrm{Ann}}
\newcommand{\Tor}{\mathrm{Tor}}
\newcommand{\Supp}{\mathrm{Supp}}\newcommand{\supp}{\mathrm{supp}}
\newcommand{\Max}{\mathrm{Max}}
\newcommand{\Coker}{\mathrm{Coker}}
\newcommand{\Stab}{\mathrm{Stab}}
\newcommand{\Irr}{\mathrm{Irr}}\newcommand{\Inf}{\mathrm{Inf}}\newcommand{\Sup}{\mathrm{Sup}}
\newcommand{\rk}{\mathrm{rk}}\newcommand{\Fil}{\mathrm{Fil}}
\newcommand{\Sim}{{\mathrm{Sim}}} \newcommand{\SL}{{\mathrm{SL}}}
\newcommand{\Spec}{{\mathrm{Spec}}} \newcommand{\SO}{{\mathrm{SO}}}
\newcommand{\SU}{{\mathrm{SU}}} \newcommand{\Sym}{{\mathrm{Sym}}}
\newcommand{\sgn}{{\mathrm{sgn}}} \newcommand{\tr}{{\mathrm{tr}}}
\newcommand{\tor}{{\mathrm{tor}}}  \newcommand{\ur}{{\mathrm{ur}}}
\newcommand{\vol}{{\mathrm{vol}}}  \newcommand{\ab}{{\mathrm{ab}}}
\newcommand{\Sh}{{\mathrm{Sh}}} \newcommand{\Ell}{{\mathrm{Ell}}}
\newcommand{\Char}{{\mathrm{Char}}}\newcommand{\Tate}{{\mathrm{Tate}}}
\newcommand{\corank}{{\mathrm{corank}}} \newcommand{\Cond}{{\mathrm{Cond}}}
\newcommand{\Inn}{{\mathrm{Inn}}} \newcommand{\Spf}{{\mathrm{Spf}}} \newcommand{\Mp}{{\mathrm{Mp}}}
 \newcommand{\Sp}{{\mathrm{Sp}}}

    \font\cyr=wncyr10  \newcommand{\Sha}{\hbox{\cyr X}}
    \newcommand{\wt}{\widetilde} \newcommand{\wh}{\widehat} \newcommand{\ck}{\check}
    \newcommand{\pp}{\frac{\partial\bar\partial}{\pi i}}
    \newcommand{\pair}[1]{\langle {#1} \rangle}
    \newcommand{\wpair}[1]{\left\{{#1}\right\}}
    \newcommand{\intn}[1]{\left( {#1} \right)}
    \newcommand{\norm}[1]{\|{#1}\|}
    \newcommand{\sfrac}[2]{\left( \frac {#1}{#2}\right)}
    \newcommand{\ds}{\displaystyle}
    \newcommand{\ov}{\overline}
    \newcommand{\Gros}{Gr\"{o}ssencharaktere}
    \newcommand{\incl}{\hookrightarrow}
    \newcommand{\lra}{\longrightarrow}
     \newcommand{\ra}{\rightarrow}
    \newcommand{\imp}{\Longrightarrow}
    \newcommand{\lto}{\longmapsto}
    \newcommand{\bs}{\backslash}
    \newcommand{\nequiv}{\equiv\hspace{-7.8pt}/}
    \theoremstyle{plain}

    %\renewcommand{\thechapter}{\Roman{chapter}}
    \newtheorem{thm}{Theorem}[section] \newtheorem{corollary}[thm]{Corollary}
    \newtheorem{lemma}[thm]{Lemma}  \newtheorem{proposition}[thm]{Proposition}
    \newtheorem {conj}[thm]{Conjecture} \newtheorem{definition}[thm]{Definition}

    %accented words
    \newcommand{\Neron}{N\'{e}ron~}
    \newcommand{\Poincare}{Poincar\'{e}~}
    \renewcommand{\theequation}{\arabic{equation}}

    \numberwithin{equation}{section}

    \renewcommand\sectionmark[1]{\markboth{\thesection. #1}}


\usepackage[top=1in, bottom=1in, left=1.25in, right=1.25in]{geometry}


\renewcommand{\labelenumii}{\alph{enumii}.}

\renewcommand{\thefootnote}{}



\begin{document}
%\author{Jingsong Chai}
\title{A note on small theta lift}
%\date{\today}
%\maketitle{}


\begin{abstract}
In this note, we use certain sesquilienar form to realize small theta lift for even orthogonal-symplectic and unitary dual pairs over p-adic fields.
\end{abstract}



	\author{Jingsong Chai}
	\address{School of Mathematics, Physics and Finance \\ Anhui Polytechnic University \\Wuhu, Anhui,  241000\\China}
	\email{jingsongchai@hotmail.com}

	\subjclass[2010]{22E50,20C33}
	\keywords{Local theta correspondence, Unitary representation, Positivity}
	%\thanks{The author is supported by a start up funding of AHPU}


	
	\maketitle
\section{Introduction}

Let $(G',G)$ be an irreducible type I reductive dual pair inside the symplectic group $Sp(F)$, where $F$ is a p-adic field. Let $\tilde{S}p(F)$ be the metaplectic two fold cover of $Sp(F)$. For any subgroup $H$ of $Sp(F)$ we let $\tilde{H}$ denote its inverse image in $\tilde{S}p(F)$. Let $(\omega, \mathcal{Y})$ be the Weil representation of $\tilde{Sp}(F)$. Let $(\pi', V_{\pi'})$ be any irreducible unitary representation of $\tilde{G}'$. Let $\mathcal{Y}^\infty, H_{\pi'}^\infty$ denote the space of smooth vectors. Then $\tilde{G}'$ acts on the algebraic tensor product $\mathcal{Y}^\infty\otimes V_{\pi'}^\infty$ by $\omega\otimes \pi'$. The space $\mathcal{Y}^\infty\otimes V_{\pi'}^\infty$ is naturally endowed with an inner product $<,>$ coming from the unitary structure of $\mathcal{Y}$ and $H_{\pi'}$. Consider the new sesquilinear form $<,>_{\pi'}$ defined by
\[
<\Phi,\Phi'>_{\pi'}=\int_{\tilde{G}'}<\Phi, (\omega\otimes \pi')(g')\Phi'>dg'
\]
for $\Phi,\Phi'\in \mathcal{Y}^\infty\otimes V_{\pi'}^\infty$. Assume $(G',G)$ and $\pi'$ are such that the above integral converges for all $\Phi,\Phi'$. Let $R$ denote the radical of $<,>_{\pi'}$. Then $<,>_{\pi'}$ factors to define a nondegenerate sesquilinear form on
\[
H(\pi'):=(\mathcal{Y}^\infty \otimes V_{\pi'}^\infty)/R.
\]
It is clear that $\tilde{G}$ acting on the first factor of $\mathcal{Y}^\infty\otimes V_{\pi'}^\infty$ preserves $<,>_{\pi'}$. Hence we have an action of $\tilde{G}$ on $V_{\pi'}$. The following conjecture is proposed by Li in \cite{L90}.  \\

$\mathbf{Conjecture}$. Suppose $<,>_{\pi'}$ is defined and nonzero. Then

(i) $<,>_{\pi'}$ is nonnegative, hence $H(\pi')$ is unitary.

(ii) $H(\pi')$ defines an irreducible unitary representation on the completion of $H(\pi')$ with respect to the inner product given by $(i)$.

(iii) The map $\pi'\to H(\pi'^\vee)$, where $\pi'^\vee$ denote the contragredient of $\pi'$, coincides with Howe's duality correspondence.  \\

Now let $(G',G)$ be either even orthogonal-symplectic or unitary dual pair. Let $\pi$ be an irreducible smooth representation of $G'$ with its contragredient $\pi^\vee$. Let's ignore the splitting of $G'\times G$ inside $\tilde{Sp}(F)$ for simplicity. Choose a nonzero element $\ell\in \Hom_{G'(W)\times G'(W)\times G(V)} (\omega\otimes \bar{\omega} \otimes \pi \otimes \pi^\vee , \BC)$. Define
\[
R:=\{ \varphi\otimes v^\vee\in \omega\otimes \pi^\vee: \ell(\varphi\otimes \varphi'\otimes v\otimes v^\vee)=0 \ \ \forall \varphi'\in \bar{\omega}, v\in \pi  \}
\]
be the radical of $\ell$. Put $H_{\ell,\psi}(\pi):= (\omega\otimes \pi^\vee)/R$.
The main result of the note is the following.
\begin{thm}
 $ H_{\ell,\psi}(\pi)$ is isomorphic to the small theta lift $\theta_\psi(\pi)$.
\end{thm}

$\bullet$ We have to restrict to even orthogonal-symplectic and unitary daul pairs since in the proof we used a result of Droschl(\cite{D23}), which is proved only for these dual pairs. Our method works in principal for other dual pairs as well once Droschl's result is available. \\

$\bullet$ In Li's conjecture, the sesquilinear form $<,>_{\pi'}$ is an example of $\ell$ in the above theorem. Thus our result can be viewed as an extension of part (ii) and (iii) in Li's conjecture to general irreducible smooth representations.  \\

\noindent \textbf{Acknowledgement:} We would like to thank Lu Hengfei for helpful discussions. This work is supported by a start up funding of AHPU and NSFC(no.12571010).
\\

\section{Proof of the main result}

Let $E$ be either $F$ or a quadratic extension of $F$. We fix a nontrivial additive character $\psi$ of F. Let $V$ and $W$ be vector spaces over $E$ equipped with nondegenerate sesquilinear forms
\[
(\cdot,\cdot):V\times V \to E \ \ \ \ \ \ \and \ \ \ \ \ \ <\cdot,\cdot>:W\times W\to E
\]
of opposite signs. Put
\[
m=\dim V  \ \ \ \ \ \ \and \ \ \ \ \ \ n=\dim W.
\]
We will consider the dual pair $(G'(W), G(V))$ arising from $V, W$, and will distinguish certain cases in the following table.

\begin{center}
\begin{tabular}{|c c c c c|} 
 \hline
 Cases &  & $G'(W)$ & $G(V)$ & $l$ \\ [0.5ex] 
 \hline
 $A$ & [E:F]=2 & $U_m$ & $U_n$ & $n-m$ \\ 
 \hline
 %$B$ & n odd & $Mp_m$ & $O_n$ & $n-m-1$ \\
 %\hline
 %$C'$ & m odd & $O_m$ & $Mp_n$ & $n-m+1$ \\
 %\hline
 $C''$ & m even & $O_m$ & $Sp_n$ & $n-m+1$ \\
 \hline
 $D$ & n even & $Sp_m$ & $O_n$ & $n-m-1$ \\ [1ex] 
 \hline
\end{tabular}
\end{center}

Let $\boldsymbol{\chi}=(\chi_V,\chi_W)$ be a pair of characters of $E^\times$ as in section 3.2 in \cite{GI14}. Associated to these datum, there is a splitting
\[
\iota=\iota_{V,W,\boldsymbol{\chi},\psi}: G(W) \times H(V) \to Mp(\mathbb{W}),
\]
where $\mathbb{W}=V\otimes_E W$.

Let $\omega_\psi$ be the Weil representation of $Mp(\mathbb{W})$. The above splitting then induces a representation, denoted as $\omega_{V,W,\boldsymbol{\chi},\psi}$, of $G(W)\times H(V)$. If $\pi$ is an irreducible smooth (genuine) representation of $G'(W)$, the maximal $\pi$-isotypic quotient of $\omega_{V,W,\boldsymbol{\chi},\psi}$ is of the form $\pi\otimes \Theta_{V,W,\boldsymbol{\chi},\psi}(\pi)$, where $\Theta_{V,W,\boldsymbol{\chi},\psi}(\pi)$ is a smooth representation of $H(V)$. We will call $\Theta_{V,W,\boldsymbol{\chi},\psi}(\pi)$ as the big theta lift of $\pi$. It is always of finite length. We use $\theta_{V,W,\boldsymbol{\chi},\psi}(\pi)$ to denote the maximal semisimple quotient of $\Theta_{V,W,\boldsymbol{\chi},\psi}(\pi)$, and call it small theta lift of $\pi$.

The following fundamental result, called Howe duality, was conjectured by Howe (\cite{How79,How89}). It was proved by Waldspuger (\cite{Wal90}) when $p\neq 2$, and in full generality by Gan-Takeda (\cite{GT16}) and Gan-Sun (\cite{GS17}).

\begin{thm}
$\theta_{V,W,\boldsymbol{\chi},\psi}(\pi)$ is either zero or irreducible. Moreover, if for irreducible smooth representations $\pi_1, \pi_2$ of $G'(W)$, we have $\theta_{V,W,\boldsymbol{\chi},\psi}(\pi_1)\cong \theta_{V,W,\boldsymbol{\chi},\psi}(\pi_2)$, then $\pi_1\cong \pi_2$.
\end{thm}

Let $W_{-}$ denote the vector space $W$ over $E$ equipped with form $-<\cdot,\cdot>$. Let $\textbf{W}=W\oplus W_{-}$. Use $\textbf{G'(W)}$ to denote the isometry group of $\textbf{W}$, and we have an embedding
\[
i: G'(W)\times G'(W) \to \textbf{G(W)}.
\]  
Note that we have 
\[
\omega_{V,\textbf{W},\boldsymbol{\chi},\psi}\circ i\cong \omega_{V,W,\boldsymbol{\chi},\psi}\otimes (\bar{\omega}_{V,W,\boldsymbol{\chi},\psi}\cdot \chi_V)
\]
as representations of $G'(W)\times G'(W)$.

Consider the following see-saw diagram

\begin{displaymath}
\xymatrix{
\textbf{G'(W)} \ar@{-}[d] \ar@{-}[rd] & G(V)\times G(V) \ar@{-}[d]\\
G'(W)\times G'(W) \ar@{-}[ru]
&\Delta G(V)}
\end{displaymath}
Then the see-saw identity says that
\begin{eqnarray}
\label{see-saw}
&\Hom_{G(V)} (\Theta_{V,W,\boldsymbol{\chi},\psi}(\pi)\otimes \Theta_{V,W,\boldsymbol{\chi},\bar{\psi}}(\pi^\vee \chi_V), \chi_W ) \nonumber  \\
&\cong \Hom_{G'(W)\times G'(W)}(\Theta_{V,\textbf{W},\boldsymbol{\chi},\psi}(\chi_W), \pi\otimes \pi^\vee \chi_V ).
\end{eqnarray}

Inside $\textbf{W}$, set
\[
W^\triangle =\{ (w,w)\in \textbf{W}:w\in W \}, \ \ \ \ W^\nabla=\{(w,-w)\in \textbf{W}: w\in W \}.
\]
Then $W^\triangle$ is a maximal isotropic subspace, and we have a complete polarization
\[
\textbf{W}=W^\triangle \oplus W^\nabla.
\]
The stabilizer of $W^\triangle$ in $\textbf{G'(W)}$ is a Siegel parabolic subgroup $\textbf{P}=M_{\textbf{P}}U_\textbf{P}$. For $s\in \bC$ and a character $\chi$ of $E^\times$, let
\[
I_{\textbf{P}}^{\textbf{G'}}(s,\chi)=\Ind_{\textbf{P}}^{\textbf{G'(W)}} (\chi |\cdot|_E^s)
\]
denote the normalized induced representation of $\textbf{G'}=\textbf{G'(W)}$, where $\chi |\cdot|^s$ is a character of $M_{\textbf{P}}$ via 
\xymatrix{
M_{\textbf{P}}=GL(W^\triangle) \ar[r]^-{\det} &E^\times \\} 
and is extended to $\textbf{P}$ with $U_{\textbf{P}}$ acting trivially.

The Weil representation $\omega_{V,\textbf{W},\boldsymbol{\chi},\psi}$ of $\textbf{G'(W)}\times G(V)$ is realized on $\CS(V\otimes W^\nabla)$, here $'\CS'$ denote the space of Schwartz function space. For $h\in H, \phi\in \CS(V\otimes W^\nabla)$, the action of $G(V)$ is given by
\[
(\omega_{V,\textbf{W},\boldsymbol{\chi},\psi}(h))\phi(x)=\chi_W(\det h)\phi(h^{-1}x).
\]
For $\textbf{g}\in \textbf{G'(W)}$, put
\[
\CF_\phi(\textbf{g})=(\omega_{V,\textbf{W},\boldsymbol{\chi},\psi}(\textbf{g})\phi)(0).
\]
One can check that $\phi\to \CF_\phi$ is a $\textbf{G'}$-equivariant map
\[
\omega_{V,\textbf{W},\boldsymbol{\chi},\psi} \to I_{\textbf{P}}^{\textbf{G'}}(-\frac{l}{2},\chi_V).
\]
Let $R_{\textbf{W},\boldsymbol{\chi},\psi}(V,\chi_W)$ be the image of this map. The following result is due to Rallis (\cite{Ra84}).

\begin{proposition}
The map $\phi\to \CF_\phi$ induces an isomorphism
\begin{displaymath}
\Theta_{V,\textbf{W},\boldsymbol{\chi},\psi}(\chi_W)\stackrel{\cong} {\longrightarrow} R_{\textbf{W},\boldsymbol{\chi},\psi}(V,\chi_W) \subset I_{\textbf{P}}^{\textbf{G'}}(-\frac{l}{2},\chi_V).
\end{displaymath}
\end{proposition}

We have the following multiplicity one result by Droschl (\cite{D23}).

\begin{proposition}
For any irreducible smooth representation $\pi$ of $G'(W)$, and any character $\chi$ of $E^\times$, we have
\[
\dim_\BC \Hom_{G'(W)\times G'(W)} (I_{\textbf{P}}^{\textbf{G}}(s,\chi), \pi\otimes \pi^\vee \chi)=1.
\]
\end{proposition}

Note that
\begin{align*}
& \Hom_{G'(W)\times G'(W)\times G(V)} (\omega_{V,W,\boldsymbol{\chi},\psi}\otimes \bar{\omega}_{V,W,\boldsymbol{\chi},\psi} \otimes \pi \chi_V^{-1}\otimes \pi^\vee , \chi_W ) \\ 
&= \Hom_{G(V)} (\Theta_{V,W,\boldsymbol{\chi},\psi}(\pi)\otimes \Theta_{V,W,\boldsymbol{\chi},\bar{\psi}}(\pi^\vee \chi_V), \chi_W )  \\
&= \Hom_{G'(W)\times G'(W)}(\Theta_{V,\textbf{W},\boldsymbol{\chi},\psi}(\chi_W), \pi\otimes \pi^\vee \chi_V ).
\end{align*}
By Rallis' result and the above multiplicity one theorem, the dimension of these Hom spaces is at most one, and it is equal to one precisely when $\Theta_{V,W,\boldsymbol{\chi},\psi}(\pi)\neq 0$ by Lemma 6.1 in \cite{GI14}. 

Take a nonzero element $\ell\in \Hom_{G'(W)\times G'(W)\times G(V)} (\omega_{V,W,\boldsymbol{\chi},\psi}\otimes \bar{\omega}_{V,W,\boldsymbol{\chi},\psi} \otimes \pi \chi_V^{-1}\otimes \pi^\vee , \chi_W ) $. Let 
\[
R:=\{ \varphi\otimes v^\vee\in \omega_{V,W,\boldsymbol{\chi},\psi}\otimes \pi^\vee: \ell(\varphi\otimes \varphi'\otimes v\otimes v^\vee)=0 \ \ \forall \varphi'\in \bar{\omega}_{V,W,\boldsymbol{\chi},\psi}, v\in \pi\chi_V^{-1}  \}
\]
be the radical of $\ell$. Put $H_{\ell,\psi}(\pi):= (\omega_{V,W,\boldsymbol{\chi},\psi}\otimes \pi^\vee)/R$.

\begin{thm}
We have 
\[
H_{\ell,\psi}(\pi)\cong \theta_{V,W,\boldsymbol{\chi},\psi}(\pi).
\]
\end{thm}
\begin{proof}
Since 
\[
\Theta_{V,W,\boldsymbol{\chi},\psi}(\pi)=(\omega_{V,W,\boldsymbol{\chi},\psi}\otimes \pi^\vee)_{G(V)},
\]
it follows that $H_{\ell,\psi}(\pi)$ is a quotient of $\Theta_{V,W,\boldsymbol{\chi},\psi}(\pi)$. Hence it is of finite length. By the Howe duality, $\Theta_{V,W,\boldsymbol{\chi},\psi}(\pi)$ has a unique quotient $\theta_{V,W,\boldsymbol{\chi},\psi}(\pi)$. Thus $\theta_{V,W,\boldsymbol{\chi},\psi}(\pi)$ is also the unique quotient of $H_{\ell,\psi}(\pi)$. Hence we can write
\[
\theta_{V,W,\boldsymbol{\chi},\psi}(\pi) = (\omega_{V,W,\boldsymbol{\chi},\psi}\otimes \pi^\vee)/R_1,
\]
where $R\subseteq R_1\subset \omega_{V,W,\boldsymbol{\chi},\psi}\otimes \pi^\vee$.

If $H_{\ell,\psi}(\pi)$ is not irreducible, then $R_1$ will contain $R$ properly. Now consider
\[
H_{\ell, \bar{\psi}}(\pi^\vee \chi_V):=(\bar{\omega}_{V,W,\boldsymbol{\chi},\psi} \otimes \pi \chi_V^{-1} )/R',
\]
where $R'\subset \bar{\omega}_{V,W,\boldsymbol{\chi},\psi} \otimes \pi \chi_V^{-1}$ is defined similarly as $R$. Also $\theta_{V,W,\boldsymbol{\chi},\bar{\psi}}(\pi^\vee \chi_V)$ is a quotient of $H_{\ell,\bar{\psi}}(\pi^\vee \chi_V)$ and can be written as 
\[
\theta_{V,W,\boldsymbol{\chi},\bar{\psi}}(\pi^\vee \chi_V)=(\bar{\omega}_{V,W,\boldsymbol{\chi},\psi} \otimes \pi \chi_V^{-1} )/R_1'
\]
with $R'\subseteq R_1' \subset \bar{\omega}_{V,W,\boldsymbol{\chi},\psi} \otimes \pi \chi_V^{-1} $.

By Corollary 18.4 in \cite{GKT}, there is a natural nonzero linear form in 
\[
\Hom_{G(V)}(\theta_{V,W,\boldsymbol{\chi},\psi}(\pi)\otimes \theta_{V,W,\boldsymbol{\chi},\bar{\psi}}(\pi^\vee \chi_V), \chi_W).
\]
View this nonzero linear form as an element in 
\[
\Hom_{G(V)}((\omega_{V,W,\boldsymbol{\chi},\psi}\otimes \pi^\vee)/R_1 \otimes (\bar{\omega}_{V,W,\boldsymbol{\chi},\psi} \otimes \pi \chi_V^{-1} )/R_1' ,\chi_W),
\]
which can be naturally extended to a nonzero liner form in 
\[
\Hom_{G(V)}( H_{\ell,\psi}(\pi)\otimes H_{\ell,\bar{\psi}}(\pi^\vee \chi_V),\chi_W ).
\]
Extend this nonzero linear form further to an element in
\[
\Hom_{G(V)}( \Theta_{V,W,\boldsymbol{\chi},\psi}(\pi)\otimes \Theta_{V,W,\boldsymbol{\chi},\bar{\psi}}(\pi^\vee \chi_V), \chi_W  ).
\]
This nonzero element will correspond to a nonzero element $\ell'$ in
\[
\Hom_{G'(W)\times G'(W)\times G(V)} (\omega_{V,W,\boldsymbol{\chi},\psi}\otimes \bar{\omega}_{V,W,\boldsymbol{\chi},\psi} \otimes \pi \chi_V^{-1}\otimes \pi^\vee , \chi_W ).
\]
As this Hom space has dimention at most one, there exists a nonzero scalar $\lambda$, such that
\[
\ell=\lambda \ell'.
\]
But, by the construction of $\ell'$, the radical is $R_1$, which is strictly larger than $R$, and we got a contradiction, which will imply that $H_{\ell,\psi}(\pi)$ is irreducible.

So $H_{\ell,\psi}(\pi)$ is an irreducible quotient of $\Theta_{V,W,\boldsymbol{\chi},\psi}(\pi)$, it has to be $\theta_{V,W,\boldsymbol{\chi},\psi}(\pi)$ by Howe duality, and the conclusion follows.
\end{proof}


%\section{Unitary preservation}








\begin{thebibliography}{SK}


\bibitem[D23]{D23}
J.Droschl. {\em Proof of a conjecture of Kudla and Rallis on quotients of degenerate principal series.}  \textit{Advances in Mathematics.} Vol.464, 110145, 2025.

\bibitem[GI14]{GI14}
Wee Tech Gan and Atsushi Ichino. {\em Formal degrees and local theta correspondence.} \textit{Inventiones Mathematicae.} Vol 195, 509-672, 2014.

\bibitem[GKT]{GKT}
Wee Tech Gan, S. Kudla and S. Takeda. {\em The local theta correspondence.} Preprint.

\bibitem[GS17]{GS17}
Wee Teck Gan and Binyong Sun. {\em The Howe duality conjecture: quaternionic case.} \textit{Representation theory, number theory, and invariant theory.} Progr. Math., 323, Birkhauser/Springer, Cham, 175-192, 2017.

\bibitem[GT16]{GT16}
Wee Teck Gan and S. Takeda. {\em A proof of the Howe duality conjecture.} \textit{J. Amer. Math. Soc.} 29, no.2, 473-493. 2016

\bibitem[How79]{How79}
Roger Howe. {\em $\theta$-series and invariant theory.} \textit{Automorphic forms, representations and L-functions.} Proc. Sympos. Pure Math.,XXXIII, Amer. Math. Soc., Providence, R.I., 1979, 275-285.

\bibitem[How89]{How89} 
Roger Howe. {\em Transcending classical invaraint theory.} \textit{J. Amer. Math. Soc.} 2, no.3, 535-552, 1989.

\bibitem[L90]{L90}
Jianshu Li. {\em Theta lifting for unitary representations with nonzero cohomology.} \textit{Duke. Math. Journal} Vol.61, no.3, 913-937, 1990


\bibitem[Wal90]{Wal90}
J.-L. Waldspurger. {\em Démonstration d’une conjecture de dualité de Howe dans le cas p-adique, $p\neq 2$.} Festschrift in honor of I.I.Piatetski-Shapiro on the occasion of his sixtieth birthday, Part I (Ramat Aviv, 1989), Israel Math. Conf. Proc., 2, Weizmann, Jerusalem, 267-324, 1990.

\bibitem[Ra84]{Ra84}
S.Rallis. {\em On the Howe duality conjecture.} \textit{Compos. Math.} 51, 339-399, 1984.

\end{thebibliography}


\end{document} 
